% ─────────────────────────────────────────────────────────────────────────────
\section{Simple Continued Fractions\texorpdfstring{\protect\footnotemark}{}}\footnotetext{As with all the other sections, the following is not meant to be an exhaustive 
treatment of the topic, but rather a sketch of the main results and algorithms related to SCFs that are relevant to \catena. 
For a more comprehensive treatment, see \cite{khinchin1964} and \cite{rockett1992}.
}
% ─────────────────────────────────────────────────────────────────────────────

In the most general sense, a continued fraction is an expression of the form
\[
  a_0 + \cfrac{b_1}{a_1 + \cfrac{b_2}{a_2 + \cfrac{b_3}{\ddots + \cfrac{b_n}{a_n} + \ldots}}}
\]
where $a_0, a_1, \ldots, a_n$ and $b_1, b_2, \ldots, b_n$, depending on the context, 
may be assumed to be integers, rational numbers, either real numbers or complex numbers,
and even functions. For the purposes of \catena, and by extension this document, we 
we restrict attention to the so-called \textit{simple continued fractions} (SCFs), that is, those 
where $b_n = 1$ for all $n$ and $a_n$ are integers with $a_1, a_2, \ldots, a_n, \ldots > 0$.\footnote{
    In the literature is common to allow $a_0 \in \R$ (for instance, see \cite{khinchin1964}), 
    but for the sake of avoiding floating-point arithmetic, as well as to gain a notion of uniqueness,
    we will restrict it to $\Z$. If nothing more, we do not lose much generality since any real number 
    admits a unique SCF expansion with integer part $a_0 = \lfloor x \rfloor$ and fractional part 
    $x - a_0 = 1/(a_1 + 1/(a_2 + \ldots))$. While we won't stop to prove it here, 
    it is worth noting that the SCF expansion of a rational number is finite, while that of an 
    irrational number is infinite.
} Namely, we are left with expressions of the form
\[
  a_0 + \cfrac{1}{a_1 + \cfrac{1}{a_2 + \cfrac{1}{\ddots + \cfrac{1}{a_n} +\cdots}}}
\]
with $a_0 \in \Z$ and $a_1, a_2, \ldots, a_n, \ldots \in \Zp$. For brevity, it is common to write this 
expression as $[a_0;\, a_1, a_2, \ldots, a_n, \ldots]$. The $a_n$ are called the \textit{partial quotients} 
of the SCF, and the $n$-th \textit{convergent} is the rational number obtained by truncating the 
expansion at $a_n$, that is, $[a_0;\, a_1, a_2, \ldots, a_n]$.

As any sequence of indexable objects, one could represent an SCF as a function $f: \N \to \Z$ such that 
$f(0) \in \Z$ and $f(n) \in \Zp$ for $n > 0$, thereby encoding the entire expansion uniformly within a 
single map. However, as suggested by the standard notation $[a_0;\, a_1, a_2, \ldots]$, it is often 
conceptually and computationally advantageous to distinguish and separate the integer part from the 
remaining partial quotients.

Accordingly, in this work we adopt a split representation: we treat the first partial quotient $a_0 \in \Z$ 
as the \textit{integer part} of the SCF, and encode the remaining coefficients via a generator $f: \N \to \Zp$ 
where $f(n) = a_{n+1}$. This distinction is not merely notational; it reflects how SCFs are 
constructed and manipulated algorithmically in \catena, and will play a central role in the 
implementations developed in the following sections.

\subsection{Core concepts and implementations}

\begin{definition}[Generating Function]
  Let $I \in \N \cup \{N\}$. A \textit{generator} indexed by $I$ is any function 
  $g: I \to \Zp$. 
  
  When $I = \N$, we say that $g$ is an \textit{infinite generator}, and when 
  $I \in \N$, we say that $g$ is a \textit{finite generator}. Furthermore, in the finite case, we can say that $g$ is 
  a I\textit{th-order} generator.
\end{definition}

\begin{remark}
    Note we allow $I = 0$. In this case, the generator is vacuous and the corresponding SCF reduces to $[a_0]$ for some $a_0 \in \Z$.
\end{remark}

\begin{catenadef}[\texttt{Generator}, \texttt{CachedGenerator}, \texttt{FiniteGenerator}]
\begin{lstlisting}[style=python]
from catena.generators import Generator, CachedGenerator, FiniteGenerator

# Infinite generator: a_n = 2 for all n > 0, with integer part a_0 = 1 (\sqrt{2} = [1; 2, 2, 2, ...])
gen = Generator(lambda n: 2)

# CachedGenerator is useful when f(n) is expensive to evaluate.
# Here, naive recursive Fibonacci is O(2^n) per call without memoisation;
# wrapping it in CachedGenerator ensures each index is computed at most once
# across all calls.
def fib(n):
    return fib(n - 1) + fib(n - 2) if n > 1 else n + 1

cgen = CachedGenerator(fib)   # fib(n) is called once per n, then cached
cgen(10) # computes fib(10) and caches results for fib(0) through fib(10)
cgen(15) # computes fib(11) through fib(15), reusing cached results as needed

# Finite generator: fixed sequence [7, 15, 1, 292, ...]  (pi = [3; 7, 15, 1, 292, ...])
fgen = FiniteGenerator([7, 15, 1, 292, 1, 1, 1, 2]) # the well-known SCF expansion of \pi begins thus
fgen(0) # 7
fgen(3) # 292
fgen(10) # raises IndexError

\end{lstlisting}
\end{catenadef}


\begin{definition}[Simple Continued Fraction]
  A \textit{simple continued fraction} is an expression of the form
  \[
    x = a_0 + \cfrac{1}{a_1 + \cfrac{1}{a_2 + \cfrac{1}{\ddots + \cfrac{1}{a_n} + \cdots}}}
  \]
  where $a_0$ is an integer and $a_1, a_2, \ldots, a_n, \ldots$ are positive integers. Thus,
  given $a_0 \in \Z$ and a generator $g: I \to \Zp$, such that $g(n) = a_{n+1}$, we represent the SCF as
  the ordered pair $x = (a_0, g)$.
\end{definition}

\begin{catenadef}[\texttt{SimpleContinuedFraction}, \texttt{FiniteSimpleContinuedFraction}]
\begin{lstlisting}[style=python]
from catena import SimpleContinuedFraction, FiniteSimpleContinuedFraction

# Infinite SCF: golden ratio phi = [1; 1, 1, 1, ...]
phi = SimpleContinuedFraction(lambda n: 1, integer_part=1)

# Finite SCF: 355/113 = [3; 7, 16]
scf = FiniteSimpleContinuedFraction.from_rational((355, 113))
scf.integer_part # 3
scf.generator    # FiniteGenerator([7, 16])
\end{lstlisting}
\end{catenadef}

\begin{definition}[Segment]
    Let $x = (a_0, g)$ be a simple continued fraction, where $g: I \to \Zp$ is a generator. We 
    define the \textit{segment} of $x$ at index $n \in I$ as the SCF obtained by truncating the first 
    $n$ partial quotients from $g$. 
    
    Formally, for $n \in I$, the segment of $x$ at index $n$ is
    given by $x_n = (a_0, g_n)$, where $g_n:=g|_n$ is the restriction of $g$ to $n$.
\end{definition}

\begin{catenadef}[\texttt{SimpleContinuedFraction.segment}]
\begin{lstlisting}[style=python]
from catena import SimpleContinuedFraction

# Example: for phi = [1; 1, 1, 1, ...], the segment at index n is also phi
phi = SimpleContinuedFraction(lambda n: 1, integer_part=1)
phi.segment(0) # [1;]
phi.segment(1) # [1; 1]
phi.segment(10) # [1; 1, 1, 1, 1, 1, 1, 1, 1, 1, 1]
\end{lstlisting}
\end{catenadef}

\begin{remark}
    Given that the segment is defined for $n \in I$, for finite SCFs, the segment is only defined 
    for $n \leq \text{size}$. For $n > \text{size}$, the segment is not defined since it would 
    require access to partial quotients that do not exist in the finite generator. Therefore, 
    in the finite case, calling \texttt{segment(n)} with $n > \text{size}$ will raise an \texttt{IndexError}.

    As for \texttt{segment(size)}, \texttt{self.shift(0)} is returned since the segment at index \texttt{size} 
    corresponds to the original SCF itself, as to avoid unnecessary copying.
\end{remark}

\begin{definition}[Remainder]
    Let $x = (a_0, g)$ be a simple continued fraction, where $g: I \to \Zp$. 
    For $n \in I$, we define the \textit{remainder} of $x$ at index $n$ as the SCF
    \[
        r_n(x) = (g(n), h),
    \]
    where $h: |I \setminus (n+1)| \to \Zp$ is defined by
    \[
        h(m) = g(m + (n+1)).
    \]
\end{definition}

\begin{remark}
    As of now, \catena does not have a dedicated implementation for remainders. We mention them 
    here because they will come up in some of the theoretical results.
\end{remark}

\begin{definition}[Convergent]
    Let $x = (a_0, g)$ be a simple continued fraction, where $g: I \to \Zp$. 
    For $n \in I$, we define the $n$-th \textit{convergent} of $x$ as the rational number
    \[
        c_n = \frac{p_n}{q_n}
    \]
    given by the value of the segment $x_n$.
\end{definition}

\begin{remark}[Segment vs Convergent]
    The segment $x_n$ is itself a (finite) SCF, while the convergent $\nicefrac{p_n}{q_n}$ 
    is a rational number. They are related by evaluation: $x_n$ represents 
    $\nicefrac{p_n}{q_n}$, and every rational number admits a unique finite SCF expansion 
    of this form.
\end{remark}

\begin{definition}[Tail]
    Let $x = (a_0, g)$ be a simple continued fraction, where $g: I \to \Zp$. For $n \in I$, the 
    \textit{tail} of $x$ is the SCF given by $t(x) = (0, g)$. This is to say, the tail of $x$ is the SCF 
    corresponding to the fractional part of $x$, $\{x\}$.

    Likewise, the n\textit{-th tail convergent} of $x$ is the $n$-th convergent of $t(x)$.
\end{definition}

\begin{catenadef}[\texttt{SimpleContinuedFraction.tail}]
\begin{lstlisting}[style=python]
from catena import SimpleContinuedFraction

# Example: for phi = [1; 1, 1, 1, ...], the tail is also phi
phi = SimpleContinuedFraction(lambda n: 1, integer_part=1)
phi.tail() # [0; 1, 1, 1, ...]

phi.tail_convergent(7) # 21/34
phi.tail().convergent(7) # 21/34
\end{lstlisting}
\end{catenadef}

\begin{theorem}[Recurrence Relations for Tail Convergents\footnote{
    See, for instance, \cite{khinchin1964}, Theorem 1, p. 4.
    For a more general version of this result, see \cite{jonesthron1980}, Theorem 2.1, p. 22.
    }]
    Let $x = (a_0, g)$ be a simple continued fraction, where $g: I \to \Zp$. 
    For $n \in I$, define the $n$-th tail convergent by
    \[
        \frac{h_n}{k_n}
        =
        [0;\, g(0), g(1), \ldots, g(n)].
    \]
    Then the numerators and denominators satisfy the recurrence relations
    \[
        h_n = g(n)\, h_{n-1} + h_{n-2}, \quad 
        k_n = g(n)\, k_{n-1} + k_{n-2} \qquad (n \geq 0),
    \]
    with initial conditions
    \[
        h_{-2} = 1,\quad h_{-1} = 0, \qquad
        k_{-2} = 0,\quad k_{-1} = 1.
    \]
\end{theorem}

\begin{remark}
  We extend the indexing of the sequences $(h_n)$ and $(k_n)$ to include 
  the values $n = -2$ and $n = -1$, with
  \[
      h_{-2} = 1,\quad h_{-1} = 0, \qquad
      k_{-2} = 0,\quad k_{-1} = 1.
  \]
  These values should be understood as formal initial conditions for the 
  recurrence, rather than as convergents. This convention allows the 
  recurrence relations to hold uniformly for all $n \geq 0$ without requiring 
  special cases, which is convenient for implementation.
\end{remark}

\begin{corollary}[Convergent from Tail Convergents]
    Let $x = (a_0, g)$ be a simple continued fraction, and let 
    $\frac{h_n}{k_n} = [0;\, g(0), \ldots, g(n)]$ be the $n$-th tail convergent. 
    Then the $n$-th convergent of $x$ is given by
    \[
        c_n = \frac{p_n}{q_n} = a_0 + \frac{h_n}{k_n}
        = \frac{a_0 k_n + h_n}{k_n}.
    \]
\end{corollary}

\begin{remark}[Caching Tail Convergents]
    The recursive nature of the tail convergents makes them ideal candidates for caching. 
    By storing previously computed tail convergents, we can avoid redundant calculations 
    and significantly improve the efficiency of algorithms that rely on them. A pretty direct
    example would be the scenario of iterating over the convergents of a given SCF: 
    by caching the tail convergents, we can compute each new convergent in constant time after 
    an initial linear time setup, rather than having to recompute the entire tail convergent 
    from scratch for each new convergent.
\end{remark}

\begin{catenadef}[\texttt{SimpleContinuedFraction.tail\_convergent}]
\begin{lstlisting}[style=python]
    from catena import SimpleContinuedFraction

    # All SCFs have a dedicated cache for tail convergents.
    phi = SimpleContinuedFraction(lambda n: 1, integer_part=1)

    phi.cache_handler # CacheHandler for tail convergents; owns the cache and handles caching logic.
    phi.tail_convergent(0) # computes and caches the 0-th tail convergent (1/1)
    phi.tail_convergent(1) # computes and caches the 1-st tail convergent (1/2), reusing the 0-th tail convergent from cache
    
    tail = phi.tail() # the tail is also phi, so it shares the same cache of tail convergents
    tail.cache_handler is phi.cache_handler # True
    tail.tail_convergent(1) # retrieves the 1-st tail convergent from cache (1/2)
\end{lstlisting}
\end{catenadef}

In the literature, tail convergents are typically not introduced explicitly, 
as they are not of primary theoretical interest. However, they provide a useful 
computational tool. In particular, they allow distinct SCFs to share both the 
same generator and the same cache of previously computed values, thereby 
avoiding redundant computations and improving efficiency.

Another point of divergence from the standard presentation is the indexing of 
convergents. In the literature, the $n$-th convergent is often defined as 
$[a_0;\, a_1, \ldots, a_n]$, which corresponds to our $(n-1)$-th convergent. 
We instead adopt zero-based indexing to reflect the conventions of Python, so 
that the first tail convergent corresponds to 
\[
    [0; g(0)] =  \frac{1}{g(0)} = \frac{1}{a_1}.
\]
While minor, this distinction is worth keeping in mind when comparing results 
with the classical literature.

\begin{remark}[RecursionError for Negative Indices]
    The recurrence relations for tail convergents are only defined for $n \geq 0$, and the seed values are defined for $n = -1$ and $n = -2$. 
    Therefore, if the recurrence is called with $n < -2$, it indicates a logic error in the implementation, as it would require access to 
    tail convergents that are not defined. In such cases, a \texttt{RecursionError} should be raised to signal that a negative value has been reached at runtime.
\end{remark}

\subsubsection{Some useful properties of convergents}

\begin{theorem}[Determinant Identity for Convergents\footnote{See, for instance, \cite{khinchin1964}, Theorem 1, p. 4.}]
    Let $x = (a_0, g)$ be a simple continued fraction, where $g: I \to \Zp$. 
    For $n \in I$, let $\frac{p_n}{q_n}$ be the $n$-th convergent of $x$. Then the following identity holds:
    \[
        p_n q_{n-1} - p_{n-1} q_n = (-1)^{n+1} \qquad (n \in I).
    \]
\end{theorem}

\begin{corollary}[Coprimality of Convergents]
    Let $x = (a_0, g)$ be a simple continued fraction, where $g: I \to \Zp$. 
    For $n \in I$, let $\frac{p_n}{q_n}$ be the $n$-th convergent of $x$. Then $p_n$ and $q_n$ are 
    coprime for all $n \in I$; that is,
    \[
        \gcd(p_n, q_n) = 1 \qquad (n \in I).
    \]
\end{corollary}

\begin{proof}
    From the determinant identity,
    \[
        p_n q_{n-1} - p_{n-1} q_n = (-1)^{n+1},
    \]
    any common divisor of $p_n$ and $q_n$ must divide the left-hand side, and hence divide $(-1)^{n+1}$. 
    Therefore, $\gcd(p_n, q_n) = 1$.
\end{proof}

\begin{remark}
    In particular, each convergent is already expressed in lowest terms. This justifies 
    the design choice of not delegating convergent arithmetic to the \texttt{Fraction} 
    class from the Python standard library, which would redundantly attempt to reduce 
    fractions.

    Since the recurrence relations preserve coprimality automatically, it is more 
    efficient to maintain numerators and denominators directly, avoiding both the 
    overhead of normalization and the cost of repeatedly instantiating \texttt{Fraction} 
    objects. Many other algorithms in \catena\ inherit this property as well, allowing 
    coprimality to be preserved without additional (and redundant) computation.
\end{remark}

\begin{corollary}[Difference of Consecutive Convergents]
    \label{coro:difference_consecutive_convergents}
    Let $x = (a_0, g)$ be a simple continued fraction, where $g: I \to \Zp$. 
    For $n \in I$, let $\frac{p_n}{q_n}$ be the $n$-th convergent of $x$. Then the following identity holds:
    \[
        \frac{p_n}{q_n} - \frac{p_{n-1}}{q_{n-1}} = \frac{(-1)^{n+1}}{q_n q_{n-1}} \qquad (n \in I).
    \]
\end{corollary}

\begin{theorem}
    Let $x = (a_0, g)$ be a simple continued fraction, where $g: I \to \Zp$. 
    For $n \in I$, let $\frac{p_n}{q_n}$ be the $n$-th convergent of $x$. Then the following identity holds:
    \[
        p_n q_{n-2} - p_{n-2} q_n = (-1)^n g(n) \qquad (n \in I, n \geq 1).
    \]
\end{theorem}

\begin{corollary}
    Let $x = (a_0, g)$ be a simple continued fraction, where $g: I \to \Zp$. 
    For $n \in I$, let $\frac{p_n}{q_n}$ be the $n$-th convergent of $x$. Then the following identity holds:
    \[
        \frac{p_n}{q_n} - \frac{p_{n-2}}{q_{n-2}} = \frac{(-1)^n g(n)}{q_n q_{n-2}} \qquad (n \in I, n \geq 1).
    \]
\end{corollary}

\begin{theorem}[Monotonicity and Convergence of Convergents]
    Let $x = (a_0, g)$ be a simple continued fraction, where $g: I \to \Zp$. 
    For $n \in I$, let $\frac{p_n}{q_n}$ be the $n$-th convergent of $x$. Then the following properties hold:
    \begin{enumerate}
        \item The sequence of convergents alternates around $x$: for even $n$, we have $c_n > x$, while for odd $n$, we have $c_n < x$.\footnote{
            Due to our indexing convention (where $c_0 = [a_0;\, a_1]$), the parity of the alternating
            inequalities is reversed relative to some standard references.
        }   
        \item The sequence of convergents is strictly decreasing for even indices and strictly increasing for odd indices.
        \item The sequence of convergents converges to $x$ as $n \to \infty$ (if $I = \N$).
    \end{enumerate}
\end{theorem}

\begin{remark}
    The monotonicity and convergence properties of convergents are fundamental to their role as approximations of the original SCF. 
    The alternating nature of the sequence ensures that each new convergent provides a better approximation than the previous one, 
    while the convergence guarantees that we can get arbitrarily close to $x$ by taking sufficiently many convergents.

    This is, among other things, the theoretical basis for being able to define a generic \texttt{\_\_float\_\_} method for SCFs, 
    which is implemented in \catena\ as the value of a fixed convergent call. Namely, \texttt{.convergent(50)} is used as the 
    default approximation for the float value of an arbitrary SCF. The \texttt{machine\_epsilon} constant from the 
    Python standard library is usually around $2^{-52} \approx 2.22 \times 10^{-16}$, so the 50-th 
    convergent of any SCF will be accurate to well within machine precision. A heuristic train of thought for this choice 
    is as follows: since the sequence of convergents alternates around $x$ and converges to $x$, we 
    focus on the value of the difference between the $n$-th convergent and $(n-1)$-th convergent, which, per Corollary \ref{coro:difference_consecutive_convergents}, is given by
    \[
        \left| \frac{p_n}{q_n} - \frac{p_{n-1}}{q_{n-1}} \right| = \frac{1}{q_n q_{n-1}}.
    \]
    On the other hand, the slowest growth rate for the denominators of convergents is given by the Fibonacci sequence, 
    which corresponds to the case where $g(n) = 1$ for all $n$. In this case, we can compute that
    \[
        \frac{1}{q_{50} q_{49}} = \frac{1}{20365011074 \cdot 32951280099} = \frac{1}{671053184118610816326} \approx 1.49 \times 10^{-21}.
    \]
    Therefore, even in the worst-case scenario, the 50-th convergent will be accurate to well within machine precision.
\end{remark}

\begin{catenadef}[\texttt{SimpleContinuedFraction.\_\_float\_\_}]
\begin{lstlisting}[style=python]
from catena import SimpleContinuedFraction

# Example: for phi = [1; 1, 1, 1, ...], the float value is approximately 1.618033988749895
phi = SimpleContinuedFraction(lambda n: 1, integer_part=1)

# We can call float() on phi directly.
float(phi) # approximately 1.618033988749895
\end{lstlisting}
\end{catenadef}

\begin{theorem}[Evaluation from Remainder\footnote{See \cite{khinchin1964}, Theorem 5, p. 7 \& Theorem 7, p. 8.}]
    Let $x = (a_0, g)$ be a simple continued fraction, where $g: I \to \Zp$. 
    For $n \in I$, let $\frac{p_n}{q_n}$ be the $n$-th convergent of $x$, and let $r_n$ be the $n$-th
    remainder of $x$. Then the following identity holds:
    \[
        x = \frac{p_{n-1}r_n + p_{n-2}}{q_{n-1}r_n + q_{n-2}} \qquad (n \in I).
    \]
\end{theorem}